\SourceLeaf{176}\par\Needspace{6\baselineskip}\medskip
\begin{center}\hypertarget{native-176-article129}{}\pdfbookmark[2]{129.}{article-129}129.\end{center}
The five methods explained so far lead directly to five more. The sole difference is that for
$x$ and $y$ we take the inclination of the orbit and the longitude of the ascending node in place of
the distances from Earth. The new methods proceed as follows.

I. From $x$ and $y$ and the two outer geocentric positions, determine the heliocentric longitudes in
the orbit and the radius vectors according to Articles 74 and 75. Use these and the corresponding
times to determine all the remaining elements, then derive from those elements the geocentric
position at the time of the middle observation. The differences of its longitude and latitude from
those of the observed position give $X$ and $Y$.

The four remaining methods share a first step. From the orientation of the orbital plane and the
geocentric positions, compute all three heliocentric longitudes in the orbit and the corresponding
radius vectors. Then proceed as follows.

II. Determine the remaining elements from the two outer positions alone and the corresponding times.
Using these elements, calculate the longitude in the orbit and the radius vector at the time of the
middle observation. The differences of these quantities from the values found earlier, that is, the
values derived from the geocentric position, give $X,Y$.

III. Alternatively, derive the remaining dimensions of the orbit from all three heliocentric
positions (Articles 82 and 83), without using the times in that calculation. Then find the time
intervals that should have elapsed in the orbit thus obtained between the first and second
observations and between the second and third. The differences of these intervals from the true ones
give $X,Y$.

IV. Calculate the remaining elements twice, first by combining the first position with the second,
then by combining the second with the third, using the corresponding time intervals. Compare the two
sets of elements and take any two of the differences for $X,Y$.

V. Finally, we may carry the same two calculations only as far as the values of the quantity denoted
by $y$ in Article 91, and then adopt the expressions given in the preceding article for $X,Y$.

The last four methods can be used safely if Earth's positions at all three observations are
sufficiently far from the nodes of the orbit. \SourceLeaf{177}The first method requires this condition to hold
only for the two outer observations, or more precisely, because the middle position may replace
either outer one, it requires that at most one of Earth's three positions lie near the nodes.
